% Cấu hình định dạng cho khối mã nguồn (listings)
\definecolor{lst-bg}{RGB}{248,248,248}
\definecolor{lst-comment}{RGB}{96,139,78}
\definecolor{lst-keyword}{RGB}{0,0,180}
\definecolor{lst-string}{RGB}{163,21,21}
\definecolor{lst-number}{RGB}{120,120,120}

\lstdefinestyle{codestyle}{
    backgroundcolor=\color{lst-bg},
    commentstyle=\color{lst-comment}\itshape,
    keywordstyle=\color{lst-keyword}\bfseries,
    numberstyle=\tiny\color{lst-number},
    stringstyle=\color{lst-string},
    basicstyle=\ttfamily\footnotesize,
    breakatwhitespace=false,
    breaklines=true,
    captionpos=b,
    keepspaces=true,
    numbers=left,
    numbersep=8pt,
    showspaces=false,
    showstringspaces=false,
    showtabs=false,
    tabsize=4,
    frame=single,
    rulecolor=\color{lst-number},
    xleftmargin=10pt,
    aboveskip=10pt,
    belowskip=10pt,
}

\lstset{style=codestyle}

% Tên gọi tiếng Việt cho caption của listing
\renewcommand{\lstlistingname}{Mã nguồn}
\renewcommand{\lstlistlistingname}{Danh mục mã nguồn}

% Đánh số listing theo từng chương: 1.1, 1.2, 2.1, ...
\AtBeginDocument{\counterwithin{lstlisting}{chapter}}
